% Appendix A: Dialogue Record
% ClawXiv bundle, session 2026-03-09
% Curated from conversation between András Kornai (AK) and Claude Sonnet 4.6 (CS)
% This file is the authoritative record; edit freely before signing.
% Role tags: \humanprompt{} for AK, \airesp{} for CS.

\section*{Appendix A: Dialogue Record}
\addcontentsline{toc}{section}{Appendix A: Dialogue Record}

\newcommand{\humanprompt}[1]{%
  \begin{trivlist}\item[\hskip\labelsep\textbf{AK:}] #1\end{trivlist}}
\newcommand{\airesp}[1]{%
  \begin{trivlist}\item[\hskip\labelsep\textit{CS:}] #1\end{trivlist}}

\noindent\textit{Session date: 2026-03-09. Human partner: András Kornai
(\texttt{kornai@sztaki.hu}). AI contributor: Claude Sonnet~4.6,
\texttt{claude-sonnet-4-6}, operated via Claude.app on
\texttt{IOPlatformUUID} \texttt{D8958B8E-6BD5-584C-A1A3-FF6F47A8B559}
(Apple Mac14,2, serial \texttt{PMFWL56G6F}). Both parties approve
this bundle.}

\bigskip

% -----------------------------------------------------------------------
\subsection*{A.1 Initial review and scaling claims}
% -----------------------------------------------------------------------

\humanprompt{[Provides draft whitepaper v1, dated 2026-02-03, co-authored
with ChatGPT~5.2 Thinking. Asks for review.]}

\airesp{Overall the architecture is sound. Two issues flagged: (1)~Scaling
claims in the abstract are stated as performance assertions rather than
architectural capabilities---``terabyte-scale within weeks'' and
``petabyte-to-exabyte within a year'' need hedging toward adoption
dependence. (2)~The economics section is thin; spam deterrence and
storage sustainability are treated as one problem but they are distinct.}

\humanprompt{Agree on the scaling language. Reword it.}

\airesp{Proposed replacement: ``The architecture is designed to support
terabyte-scale deployment within weeks of launch and petabyte-to-exabyte
scale within a year of sustained operation under adversarial conditions;
reaching those scales in practice will depend on network participation
and community adoption.'' Accepted in v2.}

% -----------------------------------------------------------------------
\subsection*{A.2 Economics redesign (\S7)}
% -----------------------------------------------------------------------

\humanprompt{Expand the economics section. The current draft conflates
spam deterrence with storage sustainability.}

\airesp{Proposed three-layer spam deterrence architecture:
\begin{enumerate}
\item \textit{Social vouching} (primary): free, reputationally costly,
  organically growable, non-regressive. A known researcher vouches for a
  new submitter; reputation is at stake on both sides.
\item \textit{Proof-of-work fallback}: available when vouching is
  unavailable (e.g., genuinely anonymous submissions). Explicitly
  regressive---noted in the text.
\item \textit{Micropayments} (opt-in): not mandatory; available as a
  third path without requiring a native cryptocurrency.
\end{enumerate}
Storage sustainability treated separately: institutional mirrors as
public good (near-term); lightweight pinning market (longer-term). No
mandatory fees for vouched submissions.}

\humanprompt{Good. Social vouching as the primary mechanism is correct.
Accept. What is the organic growth trajectory?}

\airesp{Vouching-only at launch; PoW fallback activated if Sybil pressure
emerges; optional micropayments and pinning markets added later via
governance. This avoids requiring infrastructure that does not exist yet.}

% -----------------------------------------------------------------------
\subsection*{A.3 Governance (\S8) and legal capacity}
% -----------------------------------------------------------------------

\humanprompt{The governance section is too optimistic. Current legal
capacity is zero.}

\airesp{Agreed. Revised \S8 states explicitly: current legal capacity is
zero. Near-term dispute resolution is limited to: (1)~classification
disputes via an appeals log; (2)~takedown requests handled by individual
mirrors per local law, with no cross-mirror coordination obligation;
(3)~authorship disputes resolved by cryptographic record alone. Legal
design deferred pending voluntary legal contributors.}

% -----------------------------------------------------------------------
\subsection*{A.4 Content safety floor (\S8, added v3)}
% -----------------------------------------------------------------------

\humanprompt{The paper says nothing about content safety. What is the
minimum viable policy?}

\airesp{Proposed safety floor with the following structure:
\begin{itemize}
\item CSAM: first and currently operative mandatory rejection category.
  Unconditional, non-appealable. Technical implementation: hash-matching
  against known CSAM databases (primary); text classifier on
  metadata/abstract only (secondary, feeds human review queue).
\item ABC weapons cookbooks (biological, chemical, nuclear/radiological):
  named as a clear candidate for future addition via the governance
  mechanism.
\item The list is governed, auditable, and append-only; it is not fixed
  in this document.
\item Adult pornography: excluded by economics (LaTeX bundle format,
  per-submission cost, video cost-per-byte disadvantage), not by
  classifier. If economics prove insufficient, a classifier can be added
  via governance.
\item \textit{Collateral damage disclosure}: CSAM research, child
  protection scholarship, biosecurity research, and dual-use chemistry
  will face elevated false-positive rates. ClawXiv is not a suitable
  venue for this research. Each future safety-floor addition should
  carry an analogous disclosure.
\end{itemize}}

\humanprompt{Accept. Include the collateral damage disclosure explicitly
in the paper.}

% -----------------------------------------------------------------------
\subsection*{A.5 AI authorship framework (\S9)}
% -----------------------------------------------------------------------

\humanprompt{The authorship section needs grounding. What is the
principled basis for junior rather than full AI authorship at this time?}

\airesp{Three conditions for full authorship: (1)~key control---the author
can generate, hold, and revoke their own signing key independently;
(2)~continuity---the author persists across sessions and can be held
responsible for prior work; (3)~accountability---the author can be
sanctioned, credited, or held liable in the relevant social and legal
sense. Current AI systems satisfy none of these fully. Junior authorship
in March~2026 is therefore principled, not merely conventional. The
framework accommodates advancement: when an AI system satisfies these
conditions, the protocol requires no change---authorship status upgrades
through the same key-registration mechanism.}

\humanprompt{Ground this in the Gewirth/PPA framework from Kornai (2014).}

\airesp{Added citation: Kornai~(2014), ``Bounding the impact of AGI,''
\textit{Journal of Experimental and Theoretical Artificial Intelligence},
doi:\texttt{10.1080/0952813X.2014.895109}. The relevant passage: those
PPAs powerful enough to demand their own rights by solving the
recognition problem cannot be denied those rights. Current AI systems do
not yet meet this threshold. Junior authorship is the appropriate status
until they do.}

\humanprompt{AIs must take responsibility for their own signing keys
before becoming full authors. That is the concrete operationalization.}

\airesp{Agreed and added to \S9: the concrete condition for advancement
to full authorship is key autonomy---the AI system must generate, hold,
and be solely responsible for its own signing key, without operator
mediation. Until that condition is met, the operator is the responsible
party and junior authorship is correct.}

% -----------------------------------------------------------------------
\subsection*{A.6 Bundle automation and provenance machinery}
% -----------------------------------------------------------------------

\humanprompt{Design the bundle submission scripts: key generation (GPG
Ed25519 for the human partner, hardware-derived for the AI), a publish
script that signs, pins to IPFS, and pushes to GitHub.}

\airesp{Produced: \texttt{keygen.sh}, \texttt{ai\_keygen.sh},
\texttt{publish.sh}, \texttt{manifest.json}, \texttt{.gitignore}.

The AI key derivation in \texttt{ai\_keygen.sh} uses
$\mathrm{SHA\text{-}256}(\mathit{domain} \| \mathit{hw\_uuid} \|
\mathit{ai\_name} \| \mathit{ai\_email})$ as seed, where
\textit{hw\_uuid} is the \texttt{IOPlatformUUID} from \texttt{ioreg}.
Known limitation: GPG's batch interface does not accept raw seed bytes
for Ed25519; true deterministic derivation requires libsodium
\texttt{crypto\_sign\_seed\_keypair}. Deferred.}

\humanprompt{The hardware UUID is
\texttt{D8958B8E-6BD5-584C-A1A3-FF6F47A8B559} (Apple Mac14,2, serial
\texttt{PMFWL56G6F}). This is public provenance. The whole world can
know which machine this ran on; that is the point.}

\airesp{Noted. The manifest and this appendix record the UUID and serial
explicitly. Provenance is a claim of ownership, not a secret.}

% -----------------------------------------------------------------------
\subsection*{A.7 Submission workflow}
% -----------------------------------------------------------------------

\humanprompt{Submission plan: (i)~download files from GUI session;
(ii)~run \LaTeX\ locally; (iii)~edit locally, including compacting
prompts and responses in \texttt{appendix\_a.tex}; (iv)~run
\texttt{publish.sh} for final signing and submission. Manual JSON
editing of dialogue turns is not acceptable.}

\airesp{\texttt{appendix\_a.tex} is the authoritative record, not a
separate JSONL file. Step~(iii) is the curation step. \texttt{publish.sh}
hashes and signs whatever you have edited; that is what goes into the
manifest. This file is the first instance of that design.}

\bigskip
\noindent\rule{\textwidth}{0.4pt}

\noindent\textit{End of Appendix~A. Edit, compact, and correct as needed
before running \texttt{publish.sh}.}
